\begin{abstract}
Ataques de backdoor com seleção de camadas escondem o envenenamento em alguns poucos
tensores de parâmetros, de modo que agregadores robustos que atuam tensor a tensor não
conseguem isolá-lo. Reproduzimos fielmente o ataque mais recente dessa classe,
o LaySelFL (CVPR 2026), em seu próprio cenário (ResNet18-GN, CIFAR-10, 20 clientes, 6
maliciosos, Dirichlet $\alpha{=}0{,}9$, 205 rodadas) e mostramos que ele atinge uma taxa de
sucesso do backdoor (ASR) de $82\%$ sem defesa, enquanto \emph{toda a classe} de defesas
publicadas falha em contê-lo (Multi-Krum, AlignIns, FedSurrogate, FLAME, RFLDefense $70$--$88\%$;
a melhor, LASA, $21{,}7\%$). Diagnosticamos a falha --- o envenenamento vaza pelos tensores
convolucionais, onde cada tensor isolado permanece abaixo de qualquer limiar por tensor --- e
propomos uma defesa que \emph{corrobora evidências entre camadas}: um sinal por tensor ($z$
robusto do cosseno par a par de atualizações centradas na mediana, somado à norma da atualização)
é agregado contando quantos tensores acusam cada cliente, e um cliente confirmado por pelo menos
$T$ tensores é removido. Em quatro sementes, a família por corroboração contém o ataque (expurgo
do cliente $4{,}3\pm1{,}3\%$, substituição por consenso $5{,}5\pm2{,}0\%$, corte por tensor
$10{,}5\pm6{,}3\%$), generaliza para VGG19 e EfficientNet e para o gatilho adaptativo A3FL no
CIFAR-10 ($12{,}6\%$), e custa $\leq 6{,}5\%$ de tempo extra por rodada, sem dados limpos no
servidor. Caracterizamos o modelo de ameaça sob o qual a defesa se sustenta ($\geq 2$ atacantes
em conluio) e um adversário adaptativo que conhece a defesa.
\end{abstract}
